Artificial intelligence has long assisted scientific research, but the rapid
advance of large language models and agentic scaffolds is reshaping the
landscape: a single system can now carry a whole-stage research from an initial hypothesis all the way to final published paper---a paradigm now referred to as AutoResearch. Yet existing evaluations
reveal little about how these agents operate or where they break down. Tasks
are narrowly-scoped, evaluation measures performance but not process, and
failure diagnoses lack systematic coverage or artifact-level visibility. To
address this gap, we introduce \textbf{\bench}, featuring \textbf{100
tasks} grounded in published frontier science across seven scientific domains
and the full research lifecycle: ideation, retrieval, execution, analysis,
writing, and review. Evaluating eight harness--model combinations yields
\textbf{800 autoresearch agent trajectories}, with process-level annotation.
We organize these insights into \textbf{\taxonomy} (AutoResearch Failure
Taxonomy), a framework of \textbf{45 empirically-grounded failure patterns}.
To enable scalable fine-grained attribution, we leverage a human-calibrated
agent-as-a-judge pipeline to inspect complete trajectories and intermediate
artifacts. While failure patterns span all stages of the research lifecycle,
they converge on a single overarching limitation: current agents lack a
\textbf{metacognitive loop}---the ability to check what they produced against
what they found, revise when it does not hold up, and question whether the path
they took was sound. 
The same patterns recur across all eight harness--model combinations, including the strongest models tested, locating the deficit at the model level rather than in any particular scaffold; whether orchestration-level interventions can close it is an open question this work does not test.  We publicly release \bench and \taxonomy to facilitate continued research and development in
autonomous scientific discovery.


\begin{figure}[htbp]
\centering
\includegraphics[width=0.98\textwidth]{figs/general_overview.pdf}
\caption{\textbf{Construction, rollout, and evaluation of \bench{}.}
\textbf{a,} 5,878 papers from nine domains are parsed into seven fields and
filtered to 100 tasks \add{ spanning seven domains}%
; the agent sees only the query (Premise, Tension), while
the target (KeyClaims, Conclusion) is withheld.
\textbf{b,} Each task runs once per harness--model pair as a six-stage episode
with revision, yielding 800 trajectories with all artifacts retained. The case
study traces one trajectory: six stage-wise deviations converging on a single
metacognitive deficit.
\textbf{c,} ARFT is induced bottom-up: experts annotate failures in full
trajectories, group them into patterns, and refine until all agree,
giving 45 patterns under 4 root-cause pillars. A judge agent then reviews the
full artifact set under a per-stage rubric, anchoring every issue to concrete
evidence and categorize it to an ARFT pattern, with a quality checker regenerating
weak analyses in a self-healing loop. It reaches $\kappa = 0.75$ (pattern) and
$0.83$ (taxonomy) against human labels, versus $0.53$ and $0.62$ for a
single-call LLM-as-a-judge.}
\label{fig:overview}
\end{figure}
